\documentclass[a4paper,12pt,headsepline]{report}

%----------------- PDF CONFIG ----------------- %
\pdfinfo{    
     /Title (PDF-Titel) 
     /Subject   (PDF-Thema)    
     /Author  (Jan Niklas Hollenbeck)    
} 

\title{Graphendatenbanken}
\author{Jan Niklas Hollenbeck}
\date{1.1.2000}



%----------------- PAKETE INKLUDIEREN ----------------- %

\usepackage{geometry} % Packet für Seitenrandabständex und Einstellung für Seitenränder
\usepackage[ngerman]{babel} % deutsche Silbentrennung

\usepackage{booktabs} %entzerrt die Tabellenzeilen und bietet verschieden dicke Unterteilungslinien
\usepackage{longtable} % Tabellen können sich nicht über mehrere Seiten 
\usepackage{graphicx} % kann LaTeX Grafiken einbinden


\usepackage[T1]{fontenc} % Zeichenencoding
\usepackage{lmodern} % typographische Qualität 
\frenchspacing % Schaltet den zusätzlichen Zwischenraum ab
\usepackage{fix-cm}
\usepackage{hyperref} % verwandelt alle Kapitelüberschriften, Verweise aufs Literaturverzeichnis und andere Querverweise in PDF-Hyperlinks
\usepackage{color}
\usepackage{url}

\usepackage[nottoc]{tocbibind}
\usepackage[T1]{fontenc}
\usepackage[utf8x]{inputenc}% Umlaute unter Mac werden automatisch gesetzt

% für Listings
\usepackage{listings}
\lstset{numbers=left, numberstyle=\tiny, numbersep=5pt, stepnumber=4, keywordstyle=\color{black}\bfseries\itshape, stringstyle=\ttfamily,showstringspaces=false,basicstyle=\footnotesize,captionpos=b}
\lstset{language=java}



%----------------- FARBEN DEFINIEREN ----------------- %
\definecolor{gray}{gray}{0.95} % Listingsbackground

%----------------- LAYOUT SETZEN ----------------- %
\geometry{left=2cm, right=2cm, top=2.5cm, bottom=2cm}
\linespread {1.25}\selectfont %1.25 da er von Haus aus 1.2 ist und 1,25 * 1,2 = 1,5 isch




%-------##-------##-------##------- ANFANG INHALT -------##-------##-------##-------%
\begin{document}



\pagenumbering{roman} % Seitennummer

%----------------- DECKBLATT -----------------%
 %----------------- KONFIGURATION ----------------- %
\pagestyle{empty} % enthalten keinerlei Kopf oder Fuß 
%----------------- HDA FBI Logo ----------------- %
\begin{figure}[t]
	\centering
	\includegraphics[width=0.6\textwidth]{logo_fbi.eps}
\end{figure}

%----------------- INHALT ----------------- %

\begin{center}
\Large Hochschule Darmstadt \\
\normalsize \textsc{- Fachbereich Informatik -} \\

% Whitespace
\vspace{105 pt}

\Huge Graphen Datenbanken \\ 
\normalsize
\vspace{20 pt}

Wissenschaftliches Arbeiten 2\\ 


\vspace{75 pt}


vorgelegt von \\
\vspace{5 pt}
Jan Niklas Hollenbeck \\
735992
\vspace{115 pt}

\begin{tabular}[h]{p{4cm}l l}
	Referent: & Prof. Dr. Martin Abel\\
	Korreferent: & Prof. Dr. Andreas M\"uller
\end{tabular}


\end{center}

 
%----------------- ABSTRACT -----------------%
 %----------------- KONFIGURATION ----------------- %
\pagestyle{empty} % enthalten keinerlei Kopf oder Fuß


\chapter*{Abstract} % (fold)
\label{cha:abtract}
	Lorem ipsum dolor sit amet, consetetur sadipscing elitr, sed diam nonumy eirmod tempor invidunt ut labore et dolore magna aliquyam erat, sed diam voluptua. At vero eos et accusam et justo duo dolores et ea rebum. Stet clita kasd gubergren, no sea takimata sanctus est Lorem ipsum dolor sit amet. Lorem ipsum dolor sit amet, consetetur sadipscing elitr, sed diam nonumy eirmod tempor invidunt ut labore et dolore magna aliquyam erat, sed diam voluptua. At vero eos et accusam et justo duo dolores et ea rebum. Stet clita kasd gubergren, no sea takimata sanctus est Lorem ipsum dolor sit amet.



verstand, scheinen nun doch vorueber zu Dies ist der Text sein. \\
siehe: http://janeden.net/die-praeambel
\nocite{Robinson2015}
 
 
%----------------- VERZEICHNISSE -----------------%
\tableofcontents % Inhaltverzeichnis

\pagestyle{plain} % zurueck setzen von roemische seitenanzahl

%----------------- EINLEITUNG -----------------%
\chapter{Einleitung}
längere Version des Abstracts 
\chapter{Graphendatenbanken}
\section{Graphen}
kurze erklärung zu Graphen und deren Funktionsweise.
\section{Funktionsweise Graphendatenbanken}
einführung in das Konzept der Graphendatenbanken
und abgrenzung zu anderen Datenbanken
\section{Abfrage und Analyse von Daten}
Wie werden Daten in Graphendatenbanken abgefragt und zur Wertgewinnung benutzt.
Einfaches abfragen von Komplexen Informationen

%----------------- INHALT -----------------%
\chapter{Vergleich der Graphendatenbanken}
Problem der nicht Standartisierung erläutern und die Methodik der Evaluierung und des praktischen tests veranschaulichen
\section{Standart Lösungen für Querys}
Die Vor und Nachteile von Standartlösungen bei Graphendatenbanken
\section{Funktionen der Systeme}
Wer kann was und nach welchen Kriterien wird bewertet 

%----------------- ERGEBNISSE -----------------%
\chapter{Ergebnisse}
Darstellung der Ergebnisse zu den Systemen und einzelnen Themen

\chapter{Related Work}
\nocite{Robinson2015}
\nocite{Angles2012}



%\nocite{wiki:xxx}
\nocite{Robinson2015}
\bibliographystyle{amsalpha}
\bibliography{literatur}
\pagenumbering{Roman}	
	

	
\end{document}